\documentclass[12pt]{article}
\usepackage[utf8]{inputenc}
\usepackage[T1]{fontenc}
\usepackage{hyperref}
\title{Ullatan Language Resources for K-12 Curriculum Development}
\author{Compiled by Claude Sonnet 5 (claude-sonnet-5) for Andr\'as Kornai}
\date{}
\begin{document}
\maketitle

\section{Introduction}
Ullatan (ISO 639-3 ull; also Ulladan, Katan, Kattalan, Kochuvelan) is a Southern Dravidian language of Kerala, India, in the Malayalamoid group. Search results describe it as extinct or dormant. Speakers shifted to Malayalam, and the Ulladan are a Scheduled Tribe, mainly in Kasaragod District. The searches found no primer, textbook, dictionary, grammar or audio for it.

\section{Primary Educational Resources}
None found. No school materials in Ullatan exist in the search results.

\section{Linguistic and Reference Materials}
\begin{itemize}
\item Wikipedia, Ullatan language: \url{https://en.wikipedia.org/wiki/Ullatan_language}
\item Wikipedia, Ulladan (the people): \url{https://en.wikipedia.org/wiki/Ulladan}
\item Glottolog entry for Ullatan: \url{https://glottolog.org/resource/languoid/id/ulla1237.bigmap.html}
\item ASJP word list for Ullatan (drawing on Bhattacharya 1976, per the search summary): \url{https://asjp.clld.org/languages/ULLATAN}
\end{itemize}

\section{Religious and Community Resources}
\begin{itemize}
\item Joshua Project, Ulladan in India (lists Malayalam as the primary language, not Ullatan): \url{https://joshuaproject.net/people_groups/18282}
\end{itemize}
No scripture or community organisation specific to Ullatan was found.

\section{Digital and Audio Resources}
None found.

\section{Curriculum Development Context}
The language is not spoken as a first language by anyone, and only a short word list and catalogue entries are available. A curriculum in Ullatan is not currently feasible from found sources. The practical route is a heritage or awareness component taught in Malayalam, built on the ASJP word list and on any further archival work, such as the Bhattacharya 1976 source, with the Ulladan community.

\end{document}